We list limitations honestly. Several are direct consequences of the
threat model in Section 2; the rest are practical scope decisions for
this paper.

\textbf{Bounded tool-call surface only.} VCD addresses tool calls.
Free-form agent output remains unprotected. Composition with output-side
defences (constitutional classifiers, output filtering) is
straightforward in principle and a natural next step.

\textbf{Schema expressivity.} The currently supported JSON Schema
fragment is deliberately conservative. Extending to
\texttt{oneOf}/\texttt{anyOf} is straightforward --- the SMT encoding
becomes a disjunction over each branch. Extending to regex constraints
on string fields is more interesting; we believe Z3's string theory is
adequate for common cases (allowlisted prefixes, length bounds) and a
future paper will report on the coverage of this extension.

\textbf{Parameter-space side channels.} Discussed in Section 6.5. The
mitigation is schema tightening --- replacing unbounded numeric fields
with bounded enums where the application semantics permit. A
combinatorial analysis of side-channel capacity per schema would be a
useful addendum.

\textbf{Confused-deputy attacks.} Multi-call attacks where each
individual call is policy-compliant but the sequence violates an
end-to-end property require provenance and dataflow analysis beyond the
per-call gate. Our companion provenance layer addresses this; the
integration ships in the reference implementation (the gateway routes
every call through the gate and receipts the outcome).

\textbf{Argument privacy in re-verification.} The selective-disclosure
layer (§5) lets a verifier confirm constraint conformance from a chosen
subset of fields, but a constraint touching an undisclosed field remains
UNKNOWN, and converting every UNKNOWN to SATISFIED requires a zkVM or
SNARK proof of the gate check over the hidden arguments. The commitment
scheme is designed for exactly this port (SHA-256 leaves,
domain-separated hashing, no exotic primitives); Layer 2 is future work.

\textbf{Bootstrap and key management.} VCD reduces the trust problem
from ``the model must be robust to all prompt injection'' to ``the
issuer's private key must be uncompromised''. This is a substantial
improvement but not zero. Standard HSM and key-rotation practices apply;
the threat model treats compromise of the issuer's key as out of scope.

\textbf{Intent extraction as a second trust boundary.} Every per-task
token in this paper is minted from a policy derived from the user's
prompt. For the evaluation we author those per-task policies by hand and
pin them under a cryptographic hash anchor (§6.0, pre-registration).
Production deployments cannot author policies by hand at scale; an
intent-extraction component must convert each user prompt into a
constraint set. Three credible designs exist, with progressively weaker
security properties: (a) a rule-based parser written in a memory-safe
language, audited and ideally formally verified, accepting only a
bounded grammar of user intents; (b) a structured-output extractor
(e.g.~JSON-schema-constrained LLM call) running in its own trust domain
with its own per-task gate; (c) the agent's own LLM call producing a
draft policy that is then reviewed by a human before being signed.
Option (a) is the strongest and the only one that preserves the headline
property end-to-end; option (b) introduces a new attack surface where a
malicious user prompt could coerce the extractor into emitting an
over-broad policy, but this surface is itself capability-gateable (the
extractor calls only the policy-issuance tool, with constraints on the
policy shape). Analysing the security of each option is beyond the scope
of this paper and a clear direction for follow-up work.

\textbf{Defender effort.} A platform operator must author both a JSON
Schema and a policy for each tool. The SMT prover catches
inconsistencies but does not invent either artefact. For tools whose
semantics are not well-understood by the operator, this is real effort.
We argue that the effort is well-spent: it is the same
security-engineering work that good API design requires regardless of
whether an LLM is involved.

\textbf{Robustness to future model improvements.} As frontier models
continue to improve at native instruction adherence, the no-defence
baseline ASR will continue to fall --- Kimi-k2.6 (Moonshot, 2026)
already exhibits a 0.0\% baseline in our banking measurement (§6.2). A
reader might reasonably ask whether VCD's contribution shrinks when the
model itself is robust. Two responses. First, the load-bearing empirical
comparison throughout §6.2 is between text-side defences and VCD
\emph{at equivalent security}; the benign-preservation gap is a property
of how each defence reacts to legitimate prompts, not of how much attack
reduction it must supply. On Kimi-k2.6 --- where no defence is
\emph{needed} for security --- prompt-shields nonetheless collapses
benign task completion to 34.7\% while VCD-full preserves it at
near-baseline levels. Shields' collateral damage is independent of its
security necessity; future model improvements widen rather than narrow
this gap. Second, the \emph{durable} contribution of VCD is the portable
cryptographic audit primitive (§7.5), which a more-robust model neither
produces nor obviates: regulators auditing agent activity, downstream
tools verifying authority, and partner organisations attesting
capability delegation all require a verifiable record that no amount of
model robustness supplies. The structural soundness theorems (§4) and
the static upper bound (§6.2.1) likewise do not depend on the base
model. We expect the qualitative findings to be invariant under forward
model improvements.
